% =============================================================================
% FAECO 中文期刊版式（计算机辅助设计与图形学学报模板近似）
% 字体说明：模板要求方正书宋_GBK/方正仿宋_GBK/黑体/Times New Roman。
% 方正 GBK 字库为商业字体，需作者自行购买安装；当前以系统宋体/黑体/仿宋近似。
% 安装方正字库后，将下方 FZ 字体变量替换为对应方正字体名即可。
% =============================================================================
\documentclass[twocolumn,10pt]{ctexart}

% -----------------------------------------------------------------------------
% 字体（方正 GBK 安装后替换此处）
%   正文:    \setCJKmainfont{方正书宋_GBK}
%   标题:    \setCJKsansfont{方正黑体_GBK}
%   仿宋:    \setCJKfamilyfont{fzfs}{方正仿宋_GBK}
% 当前用系统字体近似（fontset=windows: 宋体/黑体/楷体/仿宋）
% -----------------------------------------------------------------------------
\usepackage[UTF8,fontset=windows]{ctex}
\setmainfont[Path=./fonts/,Extension=.ttf,UprightFont=times,BoldFont=timesbd,ItalicFont=timesi,BoldItalicFont=timesbi]{Times New Roman}
\setCJKmainfont[Path=./fonts/,Extension=.ttf]{FZShuSong-Z01}
\setCJKsansfont[Path=./fonts/,Extension=.ttf]{SimHei}
\newCJKfontfamily\songfamily[Path=./fonts/,Extension=.ttf]{FZShuSong-Z01}
\newCJKfontfamily\fangfamily[Path=./fonts/,Extension=.ttf]{FZFangSong-Z02}
\newCJKfontfamily\heifamily[Path=./fonts/,Extension=.ttf]{SimHei}
\usepackage{amsmath,amssymb}
\usepackage{newtxmath}
\let\songti\songfamily\let\heiti\heifamily\let\fangsong\fangfamily
\usepackage{booktabs}
\usepackage{graphicx}
\usepackage{subfig}
\usepackage{float}
\usepackage{array}
\usepackage[nocompress]{cite}
\usepackage[hidelinks]{hyperref}
\usepackage{xurl}
\usepackage{algorithm}
\usepackage{algpseudocode}
\usepackage{multirow}
\usepackage{geometry}

% 版面：A4 双栏，栏宽 21.59 字符、栏间距 3 字符（以 3zw 近似）
\geometry{a4paper,left=2.0cm,right=2.0cm,top=2.5cm,bottom=2.5cm,columnsep=31.5pt}

\renewcommand{\eqref}[1]{\textup{（\ref{#1}）}}
\graphicspath{{../figures/}}
\setlength{\emergencystretch}{3em}
\Urlmuskip=0mu plus 3mu\relax

% -----------------------------------------------------------------------------
% 标题层级（模板：1级11.5pt黑体 / 2级五号黑体 / 3级10pt书宋）
% -----------------------------------------------------------------------------
\ctexset{
  section = {format=\heiti\fontsize{11.5}{15}\selectfont, beforeskip=6pt, afterskip=4pt},
  subsection = {format=\heiti\fontsize{10.5}{14}\selectfont, beforeskip=5pt, afterskip=3pt},
  subsubsection = {format=\songti\fontsize{10}{13}\selectfont, beforeskip=4pt, afterskip=2pt}
}

% 图题（书宋9pt）、表题（黑体9pt）、表注/分图题（书宋8pt）
\usepackage{caption}
\DeclareCaptionFont{songti}{\songti}
\DeclareCaptionFont{heiti}{\heiti}
\captionsetup[figure]{font={songti,small},labelfont={bf}}
\captionsetup[table]{font={heiti,small},labelfont={bf}}
\captionsetup[subfloat]{font={songti,scriptsize}}
\captionsetup{labelsep=space}

\begin{document}

\pagestyle{plain}

% -----------------------------------------------------------------------------
% 中文首部：标题/作者/单位/摘要/关键词/中图法分类号/DOI
\twocolumn[
\begin{center}
{\heiti\zihao{3} FAECO：面向预布局逻辑重综合的失效驱动候选生成引擎}\par
\vspace{6pt}
{\fangsong\zihao{4} 佟亚龙}\par
\vspace{4pt}
{\songti\fontsize{8}{12}\selectfont（国防科技大学计算机学院，长沙 410073）}\par\vspace{2pt}
{\songti\fontsize{8}{12}\selectfont E-mail：tyl2003@nudt.edu.cn（通信作者）}\par
\vspace{8pt}
\end{center}
\begin{flushleft}
{\heiti\fontsize{8.5}{13}\selectfont 摘\quad 要：}{\songti\fontsize{8.5}{13}\selectfont 工程变更命令（engineering change order，ECO）修复中，缓冲器插入与门尺寸调整仅能微调门级延迟，对逻辑级数过深的关键路径无能为力；逻辑重综合虽能改变电路结构，但其候选构造阶段将等价性验证与物理时序优化串行处理，导致“验证失败”与“收益不足”两类问题相互制约。为此，提出面向预布局逻辑重综合的失效驱动候选生成引擎 FAECO。将等价性与时序收益双目标编码为割图权重偏好，按六类失败模式自适应更新权重并重新搜索割解；将估计线载嵌入内环接受准则，通过理想线网改善阈值与简化 SPEF 复测两层门控筛选物理稳健候选。在 ISCAS89（8/8）、ITC-99（18/19）及 PicoRV32（2/3）基准上取得严格时序改善；收敛配置下平均 WNS 改善 1.07 ns，为纯门尺寸调整的 6.7 倍；效率优先模式下 ISCAS89（8/8）仅需 27 次 STA 调用即完成修复，开销削减 99.1\%。经真实布局布线验证，8 个电路中 5 个保持改善、1 个持平、2 个回退不超过 0.02 ns。FAECO 在开源工具链上实现了失效驱动逻辑重综合的闭环修复。}
\par
\vspace{4pt}
{\heiti\fontsize{8.5}{13}\selectfont 关键词：}{\songti\fontsize{8.5}{13}\selectfont 时序 ECO；重综合；双目标割；失效驱动反馈；简化 SPEF 门控；预布局逻辑重综合}\par
\vspace{2pt}
{\songti\fontsize{8.5}{13}\selectfont 中图法分类号：TP391.41}\par
\vspace{2pt}
{\songti\fontsize{8.5}{13}\selectfont DOI：10.3724/SP.J.1089.2026-00001}
\end{flushleft}
\vspace{6pt}
% 英文题录
\begin{center}
{\bfseries\fontsize{13.5}{18}\selectfont FAECO: Failure-Aware ECO Candidate Generation Engine for Pre-Placement Logic Re-Synthesis}\par
\vspace{5pt}
{\normalsize Tong Yalong\par}
\vspace{3pt}
{\small \itshape (College of Computer Science, National University of Defense Technology, Changsha 410073, China)}\par
\vspace{6pt}
\end{center}
\begin{flushleft}
{\bfseries Abstract:} {\small In timing engineering change order (ECO) repair, buffer insertion and gate sizing can only tune gate-level delays and cannot handle critical paths with excessive logic depth; logic re-synthesis can change the circuit structure, but its candidate construction stage processes equivalence checking and physical timing optimization serially, causing the mutual restriction of verification failures and insufficient timing gains. To address this, FAECO (Failure-Aware ECO), a failure-aware candidate generation engine for pre-placement timing-ECO logic re-synthesis, is proposed. It encodes both equivalence and timing preferences into cut weights, adaptively updates the weights and re-searches the cut according to six failure modes, and embeds estimated wire load into the inner-loop acceptance criterion to filter physically robust candidates via two-stage gating (an ideal-net improvement threshold and a simplified SPEF recheck). Strict timing improvements are achieved on ISCAS89 (8/8), ITC-99 (18/19), and PicoRV32 (2/3) benchmarks. Under the convergence configuration, the average WNS improvement is 1.07 ns, 6.7$\times$ that of the gate-sizing-only method; under the efficiency-oriented configuration, only 27 STA calls are needed to complete the 8/8 ISCAS89 repair, cutting the overhead by 99.1\%. After real placement and routing, 5 of the 8 circuits retain the improvement, 1 ties, and 2 regress by no more than 0.02 ns. FAECO implements a closed-loop failure-driven logic re-synthesis repair on the open-source toolchain.}\par
\vspace{3pt}
{\bfseries Key words:} {\small timing ECO; re-synthesis; bi-objective cut; failure-aware feedback; simplified SPEF gating; pre-placement logic re-synthesis}\par
\end{flushleft}
]%


% =============================================================================
% I. 引言
% =============================================================================
工程变更命令（engineering change order，ECO）按修复目标可分为功能 ECO 与时序 ECO：前者修复逻辑功能错误，后者修复时序违例。本文聚焦时序 ECO。当时序未能在综合阶段收敛时，需在已布线网表上执行局部修改。传统时序 ECO 主要依赖缓冲器插入（B）与门尺寸调整（G），通过调整单元延迟修复最差负裕量（WNS）与总负裕量（TNS）。在严重路径违例下，B/G 优化能力不足，其结构重映射与缓冲插入变体同样受候选空间限制\cite{chen2012metal,chen2012fixability,liang2012buffer,ho2010eco}。

工业实践中，将违例锥（最差负裕量路径所在的扇入逻辑锥）整体替换为时序友好补丁的补丁替换比 B/G 更稳定，但其依赖不可复现的工业数据。由于逻辑补丁构造与物理优化串行独立，补丁替换常陷入等价性验证失败与时序收益不足的互斥，收敛缓慢甚至失败。

现有 ECO 流程在候选构造阶段普遍遵循“先构造、后验证”的串行范式，面临三项未解决挑战：其一，按失败类型划分的六类反馈（F1--F6，含等价性验证失败、时序收益不足等）中，等价性验证失败（F1）与物理时序收益不足（F4）在候选生成期互斥，事后丢弃导致搜索空间爆炸；其二，逻辑修复与物理线载预估脱节，预布局阶段的理想候选在布线后大量失效；其三，开源工具链缺少将失效反馈嵌入候选生成内环的闭环引擎。

所提 FAECO 将补丁替换流程分解为三阶段流水线，在候选构造阶段同时编码等价性与时序目标，使验证失败不再只是被动丢弃，而是转化为下一轮搜索的反馈信号。主要贡献如下：

\begin{enumerate}
  \item \textbf{失效驱动的双目标割与闭环细化机制}：将等价性偏好编码为割图权重，从构造上缓解“等价性验证失败”与“时序收益不足”的互斥，以等价性验证失败、时序收益不足等六类失败（F1--F6）驱动搜索方向自适应收缩，并以简化 SPEF（标准寄生参数交换格式）门控筛选物理稳健候选。
\item \textbf{跨基准泛化}：在 ISCAS89\cite{iscas89}、ITC-99\cite{itc99} 与 PicoRV32\cite{picorv32} 异构族三类基准上验证了失效驱动候选生成流程的可迁移性。
\item \textbf{质量-效率权衡与物理收敛验证}：在 OpenROAD\cite{openroad} + SKY130 HD\cite{sky130} 真实版图验证中，8 个 ISCAS89 电路 5/8 布线后保持改善。
\end{enumerate}

% =============================================================================
% II. 相关工作
% =============================================================================
\section{相关工作}

ECO 修复主要分为三类技术路径：基于备用单元的结构重映射、基于逻辑补丁构造的功能修复，以及基于缓冲器插入与门尺寸调整的物理优化。

结构重映射以缓冲/尺寸调整与技术重映射的两阶段流程修复时序，将布线代价纳入备用单元分配\cite{ho2010eco,ho2012treco}，或以符号采样搜索修正点\cite{kravets2019symbolic}，其候选生成是“先割后验”，等价性验证失败只能事后丢弃。例如\cite{ho2010eco}的备用单元分配需预知布线后负载，在预布局阶段无法准确估计，导致等价候选在物理上常被拒绝；\cite{kravets2019symbolic}的符号采样虽扩大了搜索范围，但未将等价性成功率作为构造阶段的优化目标。

逻辑补丁构造以功能约简与逆变器图（FRAIG）表示与 SAT 证明为基础，通过切割违例锥并构造等价补丁修复功能错误\cite{jiang2010fraig,jiang2016resource}。

后续工作引入资源感知\cite{jiang2016resource}与统一功能/时序目标\cite{lin2011simult,luo2016unified}，但仍采用“先构造、后验证”的串行模式，补丁等价性由 SAT 扫荡或组合等价性检查（CEC）复核保证\cite{mishchenko2006sat,zhong2024stp}。

当等价性验证失败时，只能扩大搜索窗口或增加补丁规模重新枚举\cite{zhuo2018patch}，计算开销随锥规模指数增长，在数十万门级设计上难以收敛。

B/G 物理优化从最优缓冲插入算法\cite{liang2012buffer}发展到学习驱动的大规模尺寸调整与缓冲树生成\cite{liu2021rlsizer,huang2025phys,chen2024aito,zhang2025buffalo}，做门级微调而不改变结构，在路径级小幅度违例下有效；当违例源于逻辑级数过多或关键路径结构不合理时，只能缓解负载效应而无法压缩逻辑深度，优化存在明显天花板\cite{liu2021rlsizer,huang2025phys}。

学习驱动方法\cite{chen2024aito,zhang2025buffalo}虽提升效率，仍未脱离“不改逻辑级结构”的前提。

上述方法大多依赖商用 EDA 引擎或内部工业流程，公开基准与工艺库上的可复现实现极少；现有开源流程（Yosys\cite{yosys} + OpenSTA\cite{opensta} + OpenROAD\cite{openroad}）虽提供综合、静态时序分析（STA）与布局布线能力，但缺少将逻辑重综合与失效反馈集成的 ECO 修复引擎。据作者所知，截至投稿时，FAECO 是在开源 OSS-CAD 工具链与 SKY130 HD\cite{sky130} 公开工艺库上实现端到端失效驱动 ECO 闭环的首次工作，结论部分的实验即对该闭环的完整验证。

从应用阶段看，结构重映射与 B/G 物理优化面向后布局、后布线阶段的物理微调，逻辑补丁构造则多在综合后、布局前执行功能修复；前者仅在物理级调整时序而不改变逻辑结构，后者仅在逻辑级保证功能等价而不感知时序。FAECO 的预布局逻辑重综合位于两者之间，在候选构造阶段同时处理逻辑等价与物理预估。

FAECO 与上述三类方法的差异汇总如表~\ref{tab:related_diff}。

\begin{table}[!t]
\centering
\caption{FAECO 与已有工作的关键差异}
\label{tab:related_diff}
\footnotesize
\setlength{\tabcolsep}{2pt}
\begin{tabular}{p{1.3cm}p{6.0cm}}
\toprule
\textbf{维度} & \textbf{已有工作与 FAECO 的差异} \\
\midrule
等价性处理 & 已有工作事后验证，CEC 失败即丢弃候选；FAECO 前置为割图权重偏好，降低验证失败率，最终仍由 CEC 保证 \\
失败响应 & 已有工作扩大搜索宽度；FAECO 按失败类型自适应收缩搜索方向（如 F4 上调关键覆盖奖励） \\
物理衔接 & 已有工作逻辑修复与物理优化串行独立；FAECO 将简化 SPEF 门控嵌入内环接受准则，F6 反馈闭环 \\
候选空间 & 已有工作单一策略（B/G/R 之一）或串行组合；FAECO 四类策略混合同轮竞争，按实测接受 \\
\bottomrule
\multicolumn{2}{p{7.0cm}}{\footnotesize R/G/B/JOINT 分别指逻辑重写、门尺寸调整、缓冲器插入与多门联合枚举，定义见 \S\ref{sec:overview}。}\\
\end{tabular}
\end{table}

\section{方法}

\subsection{问题形式化}
FAECO 解决的时序 ECO 修复问题可形式化如下：给定映射网表 $G'$、目标周期 $T$ 与工艺库 $\mathcal{L}$，寻找一系列局部补丁 $p_1,\dots,p_k$，使修复后网表 $G^*$ 的 WNS 严格改善且功能与 $G'$ 等价。候选 $p$ 由一个或多个门及其边界端口构成，割集构造见 \S\ref{sec:mechanism}；符号见表~\ref{tab:symbols}。

\begin{table}[tbp]
\centering
\caption{符号表}
\label{tab:symbols}
\scriptsize
\setlength{\tabcolsep}{2pt}
\renewcommand{\arraystretch}{0.9}
\begin{tabular}{p{2.6cm}p{4.9cm}}
\toprule
\textbf{符号} & \textbf{定义} \\
\midrule
$G = (V, E)$ & 门级网图，顶点为门/端口，边为信号线 \\
$C(G, t) \subseteq V$ & 目标输出 $t$ 的扇入锥 \\
$K(G, t)$ & 由每个门的输入/输出容量边与扇入依赖边构成的 $s$-$t$ 分裂图 \\
$\mathrm{cost}(v)$ & 门 $v$ 的节点代价，见式~\eqref{eq:cost} \\
$s(v)$、$S=\sum_u s(u)^2$、$D_{\max}$ & 门 $v$ 的扇入锥规模、锥内规模平方和与最大逻辑深度 \\
$\lambda_b,\lambda_s,\lambda_v,\lambda_c,\lambda_e$ & 边界、尺寸、验证、关键覆盖与扇出等割图权重 \\
$\mathrm{score}(p)$ & 候选排序评分，见式~\eqref{eq:score} \\
F1--F6 & 六类失败类型，定义见表~\ref{tab:failures} \\
\bottomrule
\multicolumn{2}{p{7.3cm}}{\footnotesize 评分仅用于初筛，候选最终由 OpenSTA 实测接受；F1--F6 反馈作用于下一轮的割权重更新。}\\
\end{tabular}
\end{table}

\subsection{方法总览}\label{sec:overview}

面向时序 ECO 的 FAECO 对原始网表依次执行逻辑重综合、割图构建与候选生成、验证与时序分析，形成三阶段流水线，如图~\ref{fig:method_flow} 所示。候选修复策略分为四类：逻辑重写（R，基于库函数的等价布尔替换）、门尺寸调整（G）、缓冲器插入（B）与多门联合枚举（JOINT），下文沿用其缩写。各阶段如下：

\begin{enumerate}
  \item[(a)] \textbf{重综合}：将原始网表 $G$ 经 Yosys 两步综合流程\cite{yosys}映射为 SKY130 HD 标准单元网表 $G'$（Liberty 格式，记录单元逻辑功能、引脚与时序弧）\cite{liberty}，并生成等价的参考 BLIF（Berkeley Logic Interchange Format，门级逻辑网表交换格式）\cite{blif}网表，作为后续等价性回验的基准，实现细节见 \S\ref{sec:impl}。
  \item[(b)] \textbf{双目标割与候选生成}：以扇入锥 $C(G, t)$ 为基础构建 s-t 分裂图，把等价性偏好编码进割图权重：由 Liberty 布尔函数匹配计算每个门的 R 候选集，无等价替换候选的门其关键覆盖奖励置零、绑定 G/B 策略权重。该匹配只保证局部函数等价，全局等价仍由阶段 (c) 的 ABC \texttt{cec}（ABC 为开源逻辑综合与验证工具，\texttt{cec} 为其组合等价性检查命令）\cite{abc} 最终保证，故为权重偏好而非硬性保证。每轮并行生成加权全局最小割与关键路径覆盖割两类候选，前者由 Edmonds-Karp 算法与残差可达集求得，后者取违例锥与真实关键路径门集合的交集。候选按评分函数排序进入实测；出现 F1--F6 失败时按类型自适应调整割权重并重新搜索，失败类型与反馈动作见表~\ref{tab:failures}。
  \item[(c)] \textbf{简化 SPEF 门控验证}：候选级等价在生成期由结构签名保证，R 候选来自 Liberty 布尔函数匹配，G/B 按构造保持逻辑功能；接受修复后以 ABC \texttt{cec} 对全网表回验。时序分析以 OpenSTA 理想线网模式进行，改善超过阈值的候选才生成简化 SPEF 复测，复测仍严格改善才接受，否则归类为 F6 失败并反馈，细节见 \S\ref{sec:mechanism}。
\end{enumerate}

\begin{figure}[!htbp]
\centering
\includegraphics[width=0.98\columnwidth]{fig3_method_flow.png}%
\caption{FAECO 三阶段流水线：重综合 $\rightarrow$ 切割与细化 $\rightarrow$ 验证与时序分析。}
\label{fig:method_flow}
\end{figure}

\begin{algorithm}[tbp]
\footnotesize
\caption{FAECO 失效驱动修复主循环}\label{alg:main}
\begin{algorithmic}[1]
\Require 映射网表 $G'$、周期 $T$、工艺库 $\mathcal{L}$、权重初值 $\lambda^{(0)}$、分治阈值 $\theta$、最大轮数 $r_{\max}$
\Ensure 修复后网表 $G^*$、接受候选集合 $\mathcal{A}$
\State $G \gets G'$;\ \ $\mathcal{A}\gets\emptyset$;\ \ $r\gets 0$
\Repeat
  \State $r\gets r+1$;\ \ $\mathrm{cone}\gets$ \Call{ExtractCone}{$G$, 最差违例端点}
  \State \textbf{if} $|\mathrm{cone}|>\theta$ \textbf{then} $\mathrm{cone}\gets$ \Call{SplitConeByDepth}{$\mathrm{cone}$}
  \State $\mathcal{K}\gets$ \Call{BuildCutGraph}{$\mathrm{cone}$,\ $\lambda^{(r)}$};\ \ $C_0\gets$ \Call{CriticalPathCoverCut}{$\mathrm{cone}$}
  \State $\mathcal{P}\gets \{\text{R/G/B 单门候选}\}\cup \Call{JointSlidingWindow}{C_0,\ \mathrm{depth} \le 3}$
  \For{每个候选 $p$ 按 \Call{Score}{$p$} 降序（前 $w$ 个）}
    \State $\Delta\gets$ \Call{OpenSTA}{$G$ 应用 $p$ 后} $-$ \Call{OpenSTA}{$G$};\ \ $\Delta_{\mathrm{TNS}}\gets$ 对应 TNS 差
    \If{$(\Delta>0$ 或 $\Delta=0$ 且 $\Delta_{\mathrm{TNS}}>0)$ 且 \Call{VerifyEquivalence}{$p$} 通过}
      \State $G\gets G$ 加 $p$;\ \ $\mathcal{A}\gets\mathcal{A}\cup\{p\}$;\ \ \textbf{break}
    \Else
      \State 按失败类型 $f$（F1--F6）记录；$\lambda^{(r+1)}\gets$ \Call{RefineWeights}{$\lambda^{(r)}$, $f$}
    \EndIf
  \EndFor
\Until{无严格改善候选 或 $r=r_{\max}$}
\State \Return $G,\ \mathcal{A}$
\end{algorithmic}
\end{algorithm}

算法~\ref{alg:main} 中，第 3--4 行对应阶段 (a) 的锥抽取与深度分治，第 5--6 行对应阶段 (b) 的割图构建与候选生成，第 7--12 行对应阶段 (c) 的实测验证与失败反馈。候选枚举确保每轮至少一个 G/R 候选参与竞争；关键路径覆盖割抽取关键路径门集合作为割集候选，与加权全局最小割并行生成、同轮实测。

等价性预筛是生成期的结构签名检查，而非 ABC \texttt{cec} 的全网表形式验证；后者仅在接受修复后执行一次，故 CEC 开销不随候选数增长。

接受条件包含 WNS 持平但 TNS 改善的情形，与 \S\ref{sec:setup} 的收敛配置一致。每轮只接受严格改善，WNS 单调不减且以上界 0 为限，候选空间有限，故算法必然终止；20 轮是该语义下的观测上限。

\subsection{核心机制}\label{sec:mechanism}

\subsubsection{双目标割与节点代价}
\begin{figure}[!t]
\centering
\caption{双目标割与候选生成示意：扇入锥经 s-t 分裂后，加权全局最小割（红虚线）与关键路径覆盖割（黄圈）并行生成 R/G/B/JOINT 候选并进入评分排序}
\label{fig:cut_generation}
\includegraphics[width=0.9\columnwidth]{fig_cut_generation.png}
\end{figure}

在扇入锥 $C(G, t)$ 上构建 s-t 分裂图 $K(G, t)$，定义见表~\ref{tab:symbols}、示意见图~\ref{fig:cut_generation}。每个门拆成输入/输出容量边，扇入依赖用无穷容量边，节点权重 $\mathrm{cost}(v)$ 与关键路径、层级、尺寸联动；全局最小割由 Edmonds-Karp 算法（基于广度优先搜索的最短路增广最大流算法）\cite{edmondskarp} 求出，割边源端即候选门。

分裂图复杂度为 $O(VE^2)$，其中顶点数 $V=2n+2$、$n$ 为锥内门数，边数 $E=O(n+f)$、$f$ 为扇入依赖边；数十万门级设计上单次割在数秒内完成。

传统流程先按权重割、再在复核阶段注入等价检查，导致等价性验证失败（F1）与时序收益不足（F4）互斥。

FAECO 将等价性约束直接编码进割图：构建 $K(G, t)$ 时由 Liberty 布尔函数计算每个门的 R 候选集 $\mathcal{R}(v)$。对 $\mathcal{R}(v)=\emptyset$ 的不可替换门，其关键覆盖奖励置零并绑定 G/B 策略权重，使割集优先避开这类门，从而降低而非消除 F1 发生率。

有 R 候选的门保留关键覆盖奖励，并按扇出数叠加等价稳定割代价，因为高扇出门更易在等价验证中失败；G/B 与 R 在同一割集方案内共存竞争，最终取舍由逐候选 OpenSTA 实测决定。

例如 b06 关键路径上的复合门均无 R 候选，关键覆盖奖励置零后割集方案只能生成 G/JOINT 候选，实验验证见 \S\ref{sec:itc99}。

候选排序中的 $\Delta\mathrm{timing}(p)$ 定义为候选 $p$ 应用前后的 OpenSTA 实测 WNS 差，在理想线网模式、同一时序约束文件 SDC 与时钟约束下计算：
\begin{equation}\label{eq:dtiming}
\Delta\mathrm{timing}(p) = \mathrm{WNS}_{\mathrm{after}}(p) - \mathrm{WNS}_{\mathrm{before}}.
\end{equation}
生成阶段以逻辑级缩减与库延迟差作代理估计，实测阶段以真实值排序。节点代价 $\mathrm{cost}(v)$ 是一组启发式代价，其与 WNS 的关联通过关键路径覆盖奖励建立：设 $\mathrm{depth}(v)$ 为 $v$ 在锥内的逻辑深度、$\mathbf{1}_{\mathrm{cp}}(v)$ 为 $v$ 是否位于 OpenSTA 报告的关键路径上，则
{\small
\begin{equation}
\label{eq:cost}
\begin{split}
\mathrm{cost}(v) = \frac{\substack{1 + \lambda_b/(1+\mathrm{depth}(v)) + \lambda_s\,s(v)^2/S \\ +\, 0.01\lambda_v\,s(v) + 0.1\lambda_e\,\mathrm{fanout}(v)}}{1 + \lambda_c\,\mathbf{1}_{\mathrm{cp}}(v)\,\mathrm{depth}(v)/D_{\max}},
\end{split}
\end{equation}
}
其中 $\lambda_b/\lambda_s/\lambda_v/\lambda_c$ 为边界、尺寸、验证、关键覆盖权重，$\lambda_e$ 为等价稳定权重（默认均 $1.0$），$\mathrm{fanout}(v)$ 为门 $v$ 的扇出数（仅对有 R 候选的门计入），$s(v)$ 为门 $v$ 的扇入锥规模，$S=\sum_u s(u)^2$ 为平方和，$D_{\max}$ 为锥内最大逻辑深度。

归一化项 $s(v)^2/S$ 引导割集避开锥规模过大的门；验证代价项 $0.01\lambda_v s(v)$ 与 F5 反馈共同避开验证代价高的门；等价稳定项 $0.1\lambda_e\,\mathrm{fanout}(v)$ 提高高扇出门的割代价，由 F1 反馈上调 $\lambda_e$；关键覆盖奖励使关键路径门获得更低代价，将割集引向真实违例区域。

式~\eqref{eq:cost} 是割集与 WNS 之间的启发式桥梁而非精确代理，\S\ref{sec:ablation} 以 s382 敏感性实测作为代表性验证。

候选排序使用与割图节点代价作用层次不同的独立评分函数，权重记为 $\alpha_*$：
{\small
\begin{equation}\label{eq:score}
\begin{split}
\mathrm{score}(p) = {}&\alpha_t\,\Delta\mathrm{timing}(p) - \alpha_s\,|p| \\
&- \alpha_b\,|\mathrm{boundary\_in}(p)| \\
&- \alpha_b\,|\mathrm{boundary\_out}(p)| \\
&- \alpha_v\,c_{\mathrm{verify}}(p) + \alpha_e\,\mathbf{1}[\mathrm{equiv}(p)] \\
&- \alpha_{\mathrm{pen}}\,(1-\mathbf{1}[\mathrm{equiv}(p)]),
\end{split}
\end{equation}
}
其中 $\alpha_t,\alpha_s,\alpha_b,\alpha_v,\alpha_e,\alpha_{\mathrm{pen}}$ 为评分权重（默认 $1.0$，$\alpha_{\mathrm{pen}}=10.0$），与式~\eqref{eq:cost} 的割图权重作用层次不同：$\lambda_*$ 决定“割哪里”，$\alpha_*$ 只决定已生成候选的排序。

$\Delta\mathrm{timing}(p)$ 即式~\eqref{eq:dtiming} 的实测 WNS 差，$c_{\mathrm{verify}}(p)$ 为估计验证代价，$\mathbf{1}[\mathrm{equiv}(p)]$ 为结构签名等价预筛结果，为 1 表示通过。评分仅用于排序，最终接受与否由 OpenSTA 实测与 ABC \texttt{cec} 决定。

\subsubsection{失效驱动细化}
\begin{figure}[!t]
\centering
\caption{F1--F6 失效驱动反馈闭环：候选实测失败后按类型更新割权重并重割，直至接受或达到轮数上限}
\label{fig:feedback_loop}
\includegraphics[width=0.9\columnwidth]{fig_feedback_loop.png}
\end{figure}

F1--F6 失败分类驱动细化（表~\ref{tab:failures}，闭环见图~\ref{fig:feedback_loop}）。多轮细化按“割图、失败分类、权重更新、重割”循环，直至成功或达到最大迭代数；关闭反馈（权重固定）时作为消融对照。

权重更新不是简单重试，而是每次失败将对应权重加 $1.0$，默认初值均为 $1.0$：F1 提高等价稳定奖励并增加边界惩罚 $\lambda_b$，F2/F6 增加边界惩罚 $\lambda_b$，F3 增加尺寸惩罚 $\lambda_s$，F4 增加关键覆盖奖励 $\lambda_c$，F5 增加验证代价惩罚 $\lambda_v$ 并将最大锥门数减半，下限为 1；效果示例见 \S\ref{sec:ablation}。

\begin{table}[tbp]
\centering
\caption{F1--F6 失败类型与反馈动作}
\label{tab:failures}
\scriptsize
\setlength{\tabcolsep}{2pt}
\begin{tabular}{p{2.0cm}p{2.6cm}p{2.9cm}}
\toprule
\textbf{失败类型} & \textbf{检测条件} & \textbf{反馈动作} \\
\midrule
F1 等价性验证失败 & ABC \texttt{cec} 未通过 & 提高等价稳定奖励并增加边界惩罚，使此类候选在下一轮割代价中贬值 \\
F2 边界失败 & 锥边界未闭合 & 收紧锥边界，保留已映射门 \\
F3 尺寸过大 & 候选尺寸超过阈值 & 增加尺寸惩罚，强制更小割 \\
F4 时序收益不足 & $\Delta\mathrm{WNS} < \mathrm{threshold}$ & 上调关键覆盖奖励，使下一轮割集方案集中于真实关键路径门 \\
F5 验证超时 & $\mathrm{timeout} > \mathrm{threshold}$ & 增加验证代价惩罚 $\lambda_v$ 并将最大锥门数减半 \\
F6 SPEF 复测失败 & 理想线网下改善 $>$ 阈值但 SPEF 下不改善 & 标记高线载路径，增加边界惩罚以规避高扇出路径 \\
\bottomrule
\end{tabular}
\end{table}

\subsubsection{形式等价与 SPEF 门控}
候选级等价在生成期由结构签名保证，其中 R 候选来自 Liberty 布尔函数匹配，G/B 按构造保持逻辑功能；时序指标由 OpenSTA 3.1.0 在理想线网模式下实测得到。

全局等价回验采用 ABC \texttt{cec}，由 Yosys \texttt{miter -equiv} 与 \texttt{sat -prove-asserts} 实现；逐候选主导开销为 OpenSTA 实测 $0.3$--$4\,$s。

\begin{figure}[!t]
\centering
\includegraphics[width=0.95\columnwidth]{fig_spef_gate.png}
\caption{简化 SPEF 门控的两层验证流程：候选先经理想线网 OpenSTA 实测，改善超过阈值才进入简化 SPEF 复测，复测仍严格改善才接受；未通过者归类为 F6 反馈并上调边界惩罚后重割}
\label{fig:spef_gate}
\end{figure}

该门控以一次低成本复测过滤物理无效候选，将 SPEF 从外环的被动事后评估变为内环的迭代反馈信号。

\subsection{工程实现要点}\label{sec:impl}

为生成映射网表，FAECO 采用 Yosys 的 synth -noabc 与 abc -liberty 两步综合流程；工艺库与输入 Verilog 均固定提交哈希，保证流程可复现。

扇入锥 $C(G, t) = (V_C, E_C)$ 由反向广度优先搜索（BFS）抽取。对数十万门设计，按逻辑深度将违例锥切分为受门数上限限制的子锥，独立抽取候选并沿关键路径顺序生成，既避免内存与求解超限，也使候选集中于关键路径。上述设计在第~\ref{sec:exp} 节依次验证：主结果检验修复质量与效率，跨基准检验泛化性，消融与基线对比检验各机制的独立贡献，真实版图验证检验物理保持性与 SPEF 门控。

% =============================================================================
% IV. 实验
% =============================================================================
\section{实验}\label{sec:exp}

\subsection{实验设置}\label{sec:setup}

\textbf{环境与数据集}：OSS-CAD Suite（Yosys 0.67+146）、OpenSTA 3.1.0\cite{opensta}、Z3 5.0.0\cite{z3}、NetworkX 3.6.1\cite{networkx} 与 Python 3.11.9，运行于 WSL2 Ubuntu（8 核 16 线程、32 GB RAM）；固定种子与工具链版本保证流程确定性。时序 ECO 主实验采用 8 个 ISCAS89 顺序电路，经 Yosys 映射到 SKY130 HD 标准单元库，由 OpenSTA 以 0.5 ns 周期制造违例；ITC-99 与 PicoRV32 泛化实验见 \S\ref{sec:itc99}、\S\ref{sec:ood}。质量指标为 OpenSTA 报告的最差负松弛（worst negative slack，WNS）/总负松弛（total negative slack，TNS）（理想线网）与布线后 WNS，效率指标为候选 STA 调用次数；"严格改善"指 WNS 数值增大（含由负转正），下同。

\textbf{配置与参数}：全部电路统一流程，运行两种配置：效率优先（搜索宽度 $w=1$、至多 10 轮，决策层排序并在首个严格改善后即停）与收敛（20 轮至无严格改善）。JOINT 采用沿关键路径深度至多 3 的滑动窗口枚举，上限 50 组合且不含 B；扇入锥分治阈值为 1000 门。评分函数见式~\eqref{eq:score}，默认权重 $\alpha_t=\alpha_s=\alpha_b=\alpha_v=\alpha_e=1.0$、$\alpha_{\mathrm{pen}}=10.0$，作为排序权重区别于割图权重 $\lambda_*$。

\subsection{主结果：ISCAS89 修复质量与效率}\label{sec:iscas}

8 个 ISCAS89 电路全部严格改善，并通过组合等价性检查（CEC）确认功能等价（图~\ref{fig:iscas_main}）。策略分布方面：JOINT 贡献最大（s27/s382/s420/s953），G 次之，R 仅 s820，B 收益不超过 $0.05\,$ns。机理上，JOINT 沿关键路径压缩多级逻辑深度，单门 G 只调整驱动强度而不改变逻辑级数，R 依赖工艺库中的等价布尔候选、在 ISCAS89 中覆盖有限，后者与 \S\ref{sec:itc99} 中 b06 的失败同源。

\begin{figure}[tbp]
\centering
\caption{ISCAS89 顺序电路混合修复主结果（周期 0.5 ns）}
\label{fig:iscas_main}
\includegraphics[width=0.9\columnwidth]{fig_iscas89.png}
\end{figure}

\subsection{跨基准泛化验证}\label{sec:ood}

\subsubsection{ITC-99 基准族}\label{sec:itc99}

19 个可映射到 SKY130 HD 的 ITC-99 电路中 18 个 WNS 严格改善（94.7\%），8 个大规模电路 b14/b15/b17/b18/b19/b20/b21/b22 全部成功（图~\ref{fig:itc99}）。b06 是唯一失败：其关键路径由无等价 R 候选的复合门链构成，G 放大与 JOINT 枚举均因输入电容增大拖累前级，纯 G 与混合的候选全部被拒（148 次与 156 次）。策略分布以 JOINT 为主（12 个，b20/b21 收益最大 $+1.98$、$+2.16\,$ns）；19 电路合计 1693 次候选 STA 调用、约 12.5 分钟。

规模方面，b18 与 b19 分别含 37.6 万与 75.5 万门，各以 93 次与 58 次候选 STA 调用完成修复，规模增大数千倍而 STA 调用仍为数十次量级。

\begin{figure*}[t]
\centering
\caption{ITC-99 跨基准泛化：19 个可映射电路的 WNS 修复}
\label{fig:itc99}
\includegraphics[width=0.92\textwidth]{fig_itc99.png}
\end{figure*}

\begin{figure*}[t]
\centering
\includegraphics[width=0.99\textwidth]{fig_eff_combo_h.png}
\caption{PicoRV32 泛化与消融效率（a：PicoRV32 异构族跨设计泛化；b：搜索宽度 $\times$ 反馈消融；c：轮内决策与提前停止的 STA 效率）}
\label{fig:efficiency}
\end{figure*}

\subsubsection{PicoRV32 异构族}

同一 RTL 源下三个结构差异明显的子模块直接复用 ISCAS89 流程（图~\ref{fig:efficiency}a）：picorv32 与 pcpi\_mul 分别经 R、G 取得 $+1.13$、$+0.05\,$ns 改善，存储密集子模块无改善，构成与 b06 类似的结果边界。

\subsection{消融与基线对比}\label{sec:ablation}

本小节逐项剥离反馈、决策层与候选空间，检验各机制对质量与效率的贡献。

\textbf{质量}：收敛配置下，混合策略 8/8 优于纯 G 与随机基线、7/8 严格优于纯 B（图~\ref{fig:baseline}、表~\ref{tab:min_trials}），平均改善 1.07 ns 对纯 G 的 0.16 ns（约 6.7 倍）、中位数 1.13 对 0.09 ns。随机基线与混合共享 R/G 子空间，差距同时包含候选空间与排序/反馈的贡献；纯 G 数轮即达到自身收敛结果（表~\ref{tab:min_trials}），说明其瓶颈在于候选空间表达能力。

\begin{figure}[tb]
\centering
\includegraphics[width=0.95\columnwidth]{fig_baseline.png}
\caption{收敛配置竞争基线对比：左为修复后 WNS，右为候选 STA 调用次数}
\label{fig:baseline}
\end{figure}

\textbf{效率}：反馈开启将所需搜索宽度从 $\geq 4$--$8$ 降至 $\leq 2$，STA 开销节省 $44\%$--$89\%$（图~\ref{fig:efficiency}b）；叠加决策层排序与提前停止后，8 电路候选 STA 调用合计 27 次（均值 3.4 次/电路），较全候选评估的 110 次削减 75.5\%。留一法（以其他 7 电路归纳决策表、目标电路实测）8/8 成功，最终 WNS 与全数据表逐电路一致，说明决策层可跨电路复用。收敛配置的累计 STA 开销见表~\ref{tab:min_trials}（混合 8069 次、纯 G 约 3066 次、随机约 4800 次），效率优先配置的 27 次较纯 G 削减 99.1\%（较随机基线亦削减 99\% 以上）。同预算对照（s382/s820，60/500 次 STA）中，固定与随机两种非自适应顺序最终 WNS 一致，即非自适应排序不改变同预算质量，失败感知排序的价值体现在提前停止。

\textbf{策略与鲁棒性}：纯策略消融（图~\ref{fig:ablation}）表明，B 单独使用收益不超过 0.06 ns（ISCAS89 无净改善，ITC-99 中 11/19 仅小幅提升），混合收益由 R/G/JOINT 承担；单策略覆盖不完整，如 s382 纯 R 失败，而 b15、picorv32 单策略已达混合水平。参数敏感性上，SPEF 复测阈值（5/10/20 ps）与 s382 权重扰动（$\pm 100\%$）均不改变结论；热启动较冷启动节省近半 STA 开销，代表性收敛轨迹见图~\ref{fig:convergence}。

\begin{figure}[tb]
\centering
\includegraphics[width=0.95\columnwidth]{fig_ablation.png}
\caption{纯策略消融：基线、混合与纯 R/G/B 的最终 WNS（混合含 R/G/B/JOINT，窗口深度 3）}
\label{fig:ablation}
\end{figure}

\begin{figure}[b]
\centering
\caption{收敛配置下 WNS 随外环轮次收敛（每轮至多接受一个候选，数据点为轮开始时的网表 WNS）}
\label{fig:convergence}
\includegraphics[width=0.9\columnwidth]{fig_convergence.png}
\end{figure}

\begin{table}[!t]
\centering
\caption{不同策略在收敛配置（20 轮）下的最终 WNS 与 STA 开销，WNS 单位 ns}
\label{tab:min_trials}
\setlength{\tabcolsep}{2pt}
\footnotesize
\begin{tabular}{lrrrrrrrr}
\toprule
\textbf{电路} & \textbf{混合} & \textbf{混合} & \textbf{纯 G} & \textbf{纯 G} & \textbf{纯 B} & \textbf{纯 B} & \textbf{随机} & \textbf{随机} \\
 & \textbf{WNS} & \textbf{STA} & \textbf{WNS} & \textbf{STA} & \textbf{WNS} & \textbf{STA} & \textbf{WNS} & \textbf{STA} \\
\midrule
s27  & $+0.01$ & 119  & $-0.18$ & 32  & $+0.01$ & 42  & $-0.18$ & 84 \\
s382 & $+0.02$ & 910  & $-0.81$ & 447 & $-0.79$ & 700 & $-0.81$ & 578 \\
s420 & $-0.02$ & 580  & $-1.21$ & 146 & $-0.56$ & 328 & $-1.17$ & 330 \\
s641 & $0.00$  & 814  & $-1.18$ & 451 & $-0.36$ & 980 & $-1.05$ & 894 \\
s713 & $0.00$  & 784  & $-1.28$ & 313 & $-0.34$ & 974 & $-1.28$ & 699 \\
s820 & $-0.20$ & 1464 & $-1.17$ & 450 & $-0.93$ & 1740 & $-1.16$ & 635 \\
s832 & $-0.66$ & 1375 & $-1.16$ & 501 & $-1.17$ & 1236 & $-1.16$ & 640 \\
s953 & $-0.06$ & 2023 & $-1.22$ & 726 & $-1.07$ & 2782 & $-1.21$ & 944 \\
\bottomrule
\multicolumn{9}{p{0.95\columnwidth}}{\footnotesize STA 为达到各自最终 WNS 的累计候选 STA 调用次数；随机基线在 R/G 子空间内随机顺序，未启用 B/JOINT，其 WNS 与 STA 为 3 种子的均值。}\\
\end{tabular}
\end{table}

\subsection{物理验证：版图保持性与 SPEF 门控}\label{sec:phys_closure}

真实布局布线（placement and routing，P\&R）下 8 个 ISCAS89 电路中 5 个布线后保持改善、1 个持平、2 个轻微恶化不超过 0.02 ns（表~\ref{tab:pr8}），预布局修复大部分可保持到布线后。本节另设（A）独立 SPEF 负载扫描与（B）迭代式 SPEF 门控闭环两组实验，均基于同一网表。

（A）\textbf{独立 SPEF 负载扫描}：s382 的 R 候选与 b18 的 JOINT 候选的 WNS 改善随估计线长增大而衰减归零（图~\ref{fig:spef}），表明理想线网候选的物理保持边界有限。

\begin{figure}[tb]
\centering
\includegraphics[width=0.95\columnwidth]{fig_spef.png}
\caption{SPEF 负载扫描下 WNS 改善随线长的变化（s382 的 R 候选与 b18 的 JOINT 候选）}
\label{fig:spef}
\end{figure}

（B）\textbf{迭代式 SPEF 门控闭环}：s382 的 50 个与 b18 的 6 个候选在 40\,$\mu$m/unit 估计线载下全部未通过复测（图~\ref{fig:phys_gate}），说明当前线长估计偏保守、会拒绝真实版图上可能有效的候选。门控参数（复测阈值 10 ps、线长 40\,$\mu$m/unit）为经验设定的保守近似、未经系统校准，属于概念验证：其结论限定为可减少物理无效候选，最终裁决以真实 P\&R 为准；参数校准列为未来工作，见 \S\ref{sec:conclusion}。

\begin{figure}[tb]
\centering
\includegraphics[width=0.95\columnwidth]{fig_phys_gate.png}
\caption{SPEF 门控复测：s382（50 个）与 b18（6 个）候选的复测通过数（当前 40\,$\mu$m/unit 估计线载下全部未通过）}
\label{fig:phys_gate}
\end{figure}

\begin{table}[tbp]
\centering
\caption{真实 P\&R 版图验证：8 个 ISCAS89 电路布线后 WNS 对比（OpenROAD 2.0 + SKY130 HD PDK，WNS 单位 ns）}
\label{tab:pr8}
\footnotesize
\setlength{\tabcolsep}{2pt}
\begin{tabular}{lrrrrr}
\toprule
\textbf{电路} & \textbf{预布局基线} & \textbf{预布局修复} & \textbf{布线后基线} & \textbf{布线后修复} & $\Delta$ \\
\midrule
s27  & $-0.27$ & $-0.18$ & $-0.36$ & $-0.22$ & $+0.14$ \\
s382 & $-0.98$ & $-0.89$ & $-1.16$ & $-1.02$ & $+0.14$ \\
s420 & $-1.56$ & $-1.55$ & $-1.80$ & $-1.82$ & $-0.02$ \\
s641 & $-1.63$ & $-1.59$ & $-2.10$ & $-2.02$ & $+0.08$ \\
s713 & $-1.33$ & $-1.31$ & $-1.59$ & $-1.58$ & $+0.01$ \\
s820 & $-1.19$ & $-1.17$ & $-1.52$ & $-1.49$ & $+0.03$ \\
s832 & $-1.23$ & $-1.17$ & $-1.51$ & $-1.51$ & $0.00$ \\
s953 & $-1.31$ & $-1.27$ & $-1.60$ & $-1.61$ & $-0.01$ \\
\bottomrule
\multicolumn{6}{p{0.92\columnwidth}}{\footnotesize 统一映射网表与同一接受候选；s27 因面积过小采用 10\% 利用率布局规划；布线后 $\Delta$ 为修复网表相对基线网表的 WNS 变化。}\\
\end{tabular}
\end{table}

B 策略在受控保持时间场景下的补充验证见附录~\ref{app:hold}。
\section{结论}\label{sec:conclusion}

面向时序 ECO 的所提 FAECO 在候选构造阶段统一编码等价性与时序目标，并通过失效反馈与估计线载门控实现自适应搜索。ISCAS89 与 ITC-99 取得严格时序改善，PicoRV32 控制与数据密集子模块同样改善，效率优先模式大幅降低实测开销；真实版图验证表明预布局修复在布线后大部分保持。该实现基于开源 OSS-CAD 工具链与公开工艺库，在失效反馈与物理门控的端到端闭环上属首次开源实践。消融与留一法实验表明，失效反馈、决策层排序与提前停止对效率的提升可跨电路复用。收敛配置下的基线对比证实混合候选空间的优势主要来自候选空间本身：随机基线与混合共享 R/G 子空间，其差距同时包含排序与反馈的贡献。

主要局限与未来工作如下：

（1）真实版图验证目前限于 ISCAS89 基准，后续将扩展到全部基准并建立迭代式版图感知闭环。

（2）失败反馈学习与自适应预算分配在百万门级规模与更多工艺库上的效率边界尚未量化。

（3）简化 SPEF 门控的复测阈值与线长估计需按工艺系统校准，当前参数为经验近似的概念验证；对无 R 等价候选、G/JOINT 放大均拖累前级的电路如 b06，修复能力不足，需扩展工艺库与门级重组能力以扩大可修复空间。




\appendix
\section{保持时间补充数据}\label{app:hold}

作为补充场景实验，在 \texttt{set\_clock\_uncertainty -hold 0.8\,ns} 约束下测试 B 策略对最差保持松弛路径的修复：7/8 电路获得真实可测收益，s820/s832 修至满足时序、其余改善 $0.01$--$0.05\,$ns，且均未劣化建立时间 WNS。B 在保持时间场景有效（7/8 电路获得真实可测收益），而在建立时间场景收益可忽略（ISCAS89 上无净改善）；机理在于 D 端轻负载缓冲器链直接增加数据路径延迟；s27 因 D 端负载已轻、收益不足以弥补面积与寄生增量而未获净收益。

\begin{table}[H]
\centering
\caption{保持时间场景扩展：B 策略迭代修复的真实 OpenSTA 实测（\texttt{set\_clock\_uncertainty -hold 0.8\,ns}）}
\label{tab:hold}
\scriptsize
\setlength{\tabcolsep}{1pt}
\begin{tabular}{lrrrp{1.8cm}}
\toprule
\textbf{电路} & \textbf{保持基线} & \textbf{修复后} & \textbf{建立 WNS} & \textbf{修复动作} \\
\midrule
s382 & $-0.38$ & $-0.33$ & $-0.98$ & 2 轮：DFF\_2 / DFF\_13 各 1$\times$buf\_1 \\
s420 & $-0.36$ & $-0.33$ & $-1.56$ & 1 轮：DFF\_12 1$\times$buf\_1 \\
s641 & $-0.36$ & $-0.35$ & $-1.63$ & 1 轮：DFF\_4 1$\times$buf\_1 \\
s713 & $-0.36$ & $-0.35$ & $-1.33$ & 1 轮：DFF\_15 1$\times$buf\_1 \\
s820 & $-0.16$ & $+0.01$（满足时序） & $-1.19\rightarrow -1.12$ & 4 轮：DFF\_0/1/3/2 各 1$\times$buf\_1 \\
s832 & $-0.08$ & $+0.01$（满足时序） & $-1.23\rightarrow -1.16$ & 4 轮：DFF\_3/0/2/1 各 1$\times$buf\_1 \\
s953 & $-0.22$ & $-0.20$ & $-1.31$ & 2 轮：DFF\_1 / DFF\_19 各 1$\times$buf\_1 \\
s27 & $-0.36$ & $-0.36$（无改善） & $-0.27$ & ---（4 候选全被拒绝） \\
\bottomrule
\end{tabular}
\end{table}


\renewcommand{\refname}{参考文献}
\ctexset{section = {format=\heiti\fontsize{10.5}{14}\selectfont}}
\begin{thebibliography}{99}
\bibitem{ho2010eco} Ho K H, Chen Y P, Fang J W, et al. ECO timing optimization using spare cells and technology remapping[J]. IEEE Transactions on Computer-Aided Design of Integrated Circuits and Systems, 2010, 29(5): 697-710.
\bibitem{ho2012treco} Ho K H, Jiang J H R, Chang Y W. TRECO: Dynamic technology remapping for timing engineering change orders[J]. IEEE Transactions on Computer-Aided Design of Integrated Circuits and Systems, 2012, 31(11): 1723-1733.
\bibitem{kravets2019symbolic} Kravets V N, Lee N Z, Jiang J H R. Comprehensive search for ECO rectification using symbolic sampling[C]//Proceedings of the 56th Annual Design Automation Conference. New York: ACM, 2019: 1-6.
\bibitem{chen2012fixability} Chang H Y, Jiang I H R, Chang Y W. Timing ECO optimization via Bezier curve smoothing and fixability identification[J]. IEEE Transactions on Computer-Aided Design of Integrated Circuits and Systems, 2012, 31(12): 1857-1866.
\bibitem{chen2012metal} Chang H Y, Jiang I H R, Chang Y W. Timing ECO optimization using metal-configurable gate-array spare cells[C]//Proceedings of the 49th Annual Design Automation Conference. New York: ACM, 2012: 802-807.
\bibitem{jiang2010fraig} Wu B H, Yang C J, Huang C Y, et al. A robust functional ECO engine by SAT proof minimization and interpolation techniques[C]//Proceedings of IEEE/ACM International Conference on Computer-Aided Design. Piscataway: IEEE, 2010: 729-734.
\bibitem{lin2011simult} Chang H Y, Jiang I H R, Chang Y W. Simultaneous functional and timing ECO[C]//Proceedings of the 48th Annual Design Automation Conference. New York: ACM, 2011: 140-145.
\bibitem{jiang2016resource} Cheng A C, Jiang I H R, Jou J Y. Resource-aware functional ECO patch generation[C]//Proceedings of the Design, Automation \& Test in Europe Conference. Piscataway: IEEE, 2016: 1036-1041.
\bibitem{luo2016unified} Hung J H, Lin Y C, Cheng W K, et al. Unified approach for simultaneous functional and timing ECO[J]. IET Circuits, Devices \& Systems, 2016, 10(6): 514-521.
\bibitem{zhuo2018patch} Dao A Q, Lee N Z, Chen L C, et al. Efficient computation of ECO patch functions[C]//Proceedings of the 55th Annual Design Automation Conference. New York: ACM, 2018: 1-6.
\bibitem{mishchenko2006sat} Zhu Q, Kitchen N, Kuehlmann A, et al. SAT sweeping with local observability don't-cares[C]//Proceedings of the 43rd Annual Design Automation Conference. New York: ACM, 2006: 229-234.
\bibitem{zhong2024stp} Pan H, Zhang R, Xia Y, et al. A semi-tensor product based circuit simulation for SAT-sweeping[C]//Proceedings of the Design, Automation \& Test in Europe Conference. Piscataway: IEEE, 2024: 1-6.
\bibitem{liang2012buffer} Li Z, Zhou Y, Shi W. O(mn) time algorithm for optimal buffer insertion of nets with m sinks[J]. IEEE Transactions on Computer-Aided Design of Integrated Circuits and Systems, 2012, 31(3): 437-441.
\bibitem{liu2021rlsizer} Lu Y C, Nath S, Khandelwal V, et al. RL-Sizer: VLSI gate sizing for timing optimization using deep reinforcement learning[C]//Proceedings of the 58th Annual Design Automation Conference. New York: ACM, 2021: 733-738.
\bibitem{huang2025phys} Ye Y, Xu P, Ren L, et al. Learning-driven physically aware large-scale circuit gate sizing[J]. IEEE Transactions on Computer-Aided Design of Integrated Circuits and Systems, 2025, 44(5): 1901-1914.
\bibitem{chen2024aito} Wu H, Huang Z, Li X, et al. AiTO: Simultaneous gate sizing and buffer insertion for timing optimization with GNNs and RL[J]. Integration, 2024, 98: 102211.
\bibitem{zhang2025buffalo} Hsiao H H, Lu Y C, Lim S K, et al. BUFFALO: PPA-configurable LLM-based buffer tree generation[C]//Proceedings of IEEE/ACM International Conference on Computer-Aided Design. Piscataway: IEEE, 2025: 1-9.
\bibitem{yosys} Wolf C. Yosys open synthesis suite[EB/OL]. YosysHQ, 2024. https://yosyshq.net/yosys/.
\bibitem{liberty} Synopsys Inc. Liberty User Guide, Version 2016.12[Z]. Mountain View: Synopsys, 2016.
\bibitem{blif} Berkeley Logic Interchange Format (BLIF)[R]. Berkeley: University of California, UCB/ERL M92/126, 1992.
\bibitem{abc} Brayton R, Mishchenko A. ABC: An academic industrial-strength verification tool[C]//Proceedings of the 22nd International Conference on Computer Aided Verification (CAV). Berlin: Springer, 2010: 24-40.
\bibitem{edmondskarp} Edmonds J, Karp R M. Theoretical improvements in algorithmic efficiency for network flow problems[J]. Journal of the ACM, 1972, 19(2): 248-264.
\bibitem{iscas89} Brglez F, Bryan D, Kozminski K. Combinational profiles of sequential benchmark circuits[C]//Proceedings of IEEE International Symposium on Circuits and Systems. Piscataway: IEEE, 1989: 1929-1934.
\bibitem{itc99} Davidson S. ITC'99 benchmark circuits: preliminary results[C]//Proceedings of IEEE International Test Conference. Piscataway: IEEE, 1999: 1125-1125.
\bibitem{picorv32} Wolf C. PicoRV32: A size-optimized RISC-V CPU[EB/OL]. 2024. https://github.com/YosysHQ/picorv32.
\bibitem{openroad} Ajayi T, Chhabria V A, Fogaça M, et al. Toward an open-source digital flow: First learnings from the OpenROAD project[C]//Proceedings of the 56th ACM/IEEE Design Automation Conference. New York: ACM, 2019: 1-4.
\bibitem{sky130} SkyWater Technology Foundry. SkyWater SKY130 PDK[EB/OL]. 2024. https://github.com/google/skywater-pdk.
\bibitem{opensta} Parallax Software Inc. OpenSTA: Gate-level static timing verifier[EB/OL]. 2024. https://github.com/parallaxsw/OpenSTA.
\bibitem{z3} de Moura L, Bjørner N. Z3: An efficient SMT solver[C]//Proceedings of the 14th International Conference on Tools and Algorithms for the Construction and Analysis of Systems. Berlin: Springer, 2008: 337-340.
\bibitem{networkx} Hagberg A A, Schult D A, Swart P J. Exploring network structure, dynamics, and function using NetworkX[C]//Proceedings of the 7th Python in Science Conference. Pasadena: SciPy, 2008: 11-15.
\end{thebibliography}

\end{document}
